% Part:sets-functions-relations
% Chapter: size-of-sets
% Section: enumerations

\documentclass[../../../include/open-logic-section]{subfiles}

\begin{document}

\olfileid{sfr}{siz}{enm}

\olsection{Enumerasi lan Himpunan kang \usetoken{S}{enumerable}}

\begin{editorial}
  Perangan iki ngrembug enumerasi himpunan, kang ditetepake minangka
  surjeksi saka $\PosInt$. Andharane diwenehake alon-alon kanggo pamaca
  kang durung akeh sinau matematika. Versi alternatif kang luwih ringkes
  ana ing \olref[enm-alt]{sec}, kang netepake enumerasi
  kanthi cara beda: minangka bijeksi karo $\Nat$ utawa perangan wiwitane.
\end{editorial}

\begin{explain}
Kita wis menehi tuladha himpunan kanthi ndhaptar !!{element}s.
Saiki ayo dirembug kanthi luwih umum kepriye lan kapan kita bisa
ndhaptar !!{element}s sawijining himpunan, sanajan himpunan iku tanpa wates.
\end{explain}

\begin{defn}[Enumerasi, kanthi informal]
Kanthi informal, \emph{enumerasi} sawijining himpunan~$A$ yaiku dhaptar
(kang bisa tanpa wates) !!{element}s saka~$A$ saengga saben !!{element} saka $A$
katon ing dhaptar ing sawijining posisi winates. Yen $A$ duwe
enumerasi, mula $A$ diarani \emph{!!{enumerable}}.
\end{defn}

\begin{explain}
Sawetara bab ngenani enumerasi:
\begin{enumerate}
\item Mung dhaptar kang duwe wiwitan lan saben !!{element}
  saliyane kang kapisan duwe siji !!{element} kang langsung
  ndhisiki kang kita anggep enumerasi. Kanthi tembung liya,
  mung ana anggota kanthi cacah winates ing antarane !!{element} kapisan
  lan saben !!{element} liyane. Mligine, iki ateges saben
  !!{element} enumerasi duwe posisi winates: !!{element} kapisan
  ing posisi~$1$, kapindho ing posisi~$2$, lan sateruse.
\item Himpunan~$A$ kang padha bisa duwe enumerasi beda
  miturut urutan muncule !!{element}s: $4$, $1$,
  $25$, $16$,~$9$ padha-padha ngenumerasi himpunan limang
  wilangan kuadrat kapisan kaya $1$, $4$, $9$, $16$,~$25$.
\item Enumerasi kang mbaleni anggota tetep enumerasi: $1$, $1$, $2$,
  $2$, $3$, $3$,~\dots{} ngenumerasi himpunan kang padha karo $1$, $2$,
  $3$,~\dots{}.
\item Urutan lan pangulangan \emph{pancen} wigati nalika kita nemtokake
  sawijining enumerasi: wilangan wutuh positif bisa dienumerasi wiwit
  $1$, $2$, $3$, $1$, \dots{}, nanging polane luwih gampang katon
  yen dienumerasi kanthi cara lumrah $1$, $2$, $3$, $4$,~\dots
\item Enumerasi kudu duwe wiwitan: \dots, $3$, $2$, $1$ dudu
  enumerasi wilangan wutuh positif amarga ora duwe !!{element}
  kapisan. Kanggo mangerteni iki saka definisi informal,
  takona: "wilangan 76 katon ing posisi pira ing dhaptar?"
\item Iki dudu enumerasi wilangan wutuh positif:
  $1$, $3$, $5$, \dots, $2$, $4$, $6$, \dots\@ Masalahe yaiku
  wilangan genep ana ing posisi $\infty + 1$, $\infty + 2$, $\infty +
  3$, dudu ing posisi winates.
\item Himpunan kosong bisa didhaptar: enumerasine yaiku dhaptar kosong!{}
\end{enumerate}
\end{explain}

\begin{prop}
  Yen $A$ duwe enumerasi, mula duwe enumerasi tanpa
  pangulangan.
\end{prop}

\begin{proof}
  Upamakna $A$ duwe enumerasi $x_1$, $x_2$, \dots{} kang saben
  $x_i$ iku !!{element} saka~$A$. Pangulangan bisa dibuwang saka
  enumerasi kanthi mbuwang !!{element}s kang dibaleni. Upamane,
  enumerasi mau bisa diowahi dadi dhaptar anyar kang ngemot $x_i$ yen
  iku !!a{element} saka~$A$ kang durung ana ing antarane $x_1$, \dots,
  $x_{i-1}$, utawa mbuwang $x_i$ saka dhaptar yen wis ana ing antarane
  $x_1$, \dots,~$x_{i-1}$.
\end{proof}

Argumen pungkasan nuduhake yen kanggo mangerteni
enumerasi lan himpunan kang !!{enumerable} kanthi becik sarta mbuktekake sipate,
kita mbutuhake definisi kang luwih tliti. Definisi ing ngisor iki nyukupi kabutuhan iku.

\begin{defn}[Enumerasi, kanthi formal]
\emph{Enumerasi} sawijining himpunan $A \neq \emptyset$ yaiku sembarang
fungsi !!{surjective} $f \colon \PosInt \to A$.
\end{defn}

\begin{explain}
Ayo kita priksa yen definisi formal lan definisi informal kang nggunakake
dhaptar kang bisa tanpa wates iku ekuivalen. Kapisan, saben fungsi
!!{surjective} saka $\PosInt$ menyang himpunan~$A$ ngenumerasi~$A$.
Fungsi iku nemtokake enumerasi miturut definisi informal ing ndhuwur:
dhaptar $f(1)$, $f(2)$, $f(3)$, \dots. Amarga $f$ iku !!{surjective},
saben !!{element} saka~$A$ mesthi dadi nilai~$f(n)$ kanggo
sawijining~$n \in \PosInt$. Mula, saben !!{element} saka $A$ katon
ing posisi winates ing dhaptar. Amarga fungsi iku bisa wae ora
!!{injective}, dhaptare bisa ngemot pangulangan, nanging iki ditampa,
kaya kang wis diterangake ing ndhuwur.

Kosok baline, yen diwenehi dhaptar kang ngenumerasi kabeh !!{element}s
saka~$A$, kita bisa netepake fungsi !!a{surjective} $f\colon \PosInt \to A$
kanthi netepake $f(n)$ minangka !!{element} kaping $n$ ing dhaptar,
utawa !!{element} pungkasan yen !!{element} kaping $n$ ora ana.
Siji-sijine kasus nalika cara iki ora ngasilake fungsi !!a{surjective}
yaiku nalika $A$ kosong, mula dhaptare uga kosong. Dadi, saben dhaptar
kang ora kosong nemtokake fungsi !!a{surjective} $f\colon \PosInt \to A$.
\end{explain}

\begin{defn}
  \ollabel{defn:enumerable}
  Himpunan~$A$ iku !!{enumerable} yen lan mung yen kosong utawa duwe enumerasi.
\end{defn}

\begin{ex}
Fungsi kang ngenumerasi wilangan wutuh positif ($\PosInt$) yaiku
fungsi identitas kang diwenehake dening $f(n) = n$. Fungsi kang ngenumerasi
wilangan natural $\Nat$ yaiku fungsi $g(n) = n - 1$.
\end{ex}

\begin{ex}
Fungsi $f\colon \PosInt \to \PosInt$ lan $g \colon \PosInt \to
\PosInt$ kang diwenehake dening
\begin{align*}
f(n) & = 2n \text{ lan}\\
g(n) & = 2n - 1
\end{align*}
ngenumerasi wilangan wutuh positif genep lan wilangan wutuh positif ganjil,
kanthi runtut. Nanging loro-lorone dudu enumerasi
$\PosInt$, amarga loro-lorone ora !!{surjective}.
\end{ex}

\begin{prob}
Tetepna enumerasi wilangan kuadrat positif $1$, $4$, $9$, $16$, \dots
\end{prob}

\begin{ex}
Fungsi $f(n) = (-1)^{n} \lceil \frac{(n-1)}{2}\rceil$ (kanthi
$\lceil x \rceil$ nuduhake fungsi \emph{pambunderan munggah}, kang mbunderake
$x$ munggah menyang wilangan wutuh paling cedhak) ngenumerasi himpunan
wilangan wutuh~$\Int$. Gatekna kepriye $f$ ngasilake nilai-nilai $\Int$
kanthi "mlumpat" bolak-balik antarane wilangan wutuh positif lan negatif:
\[
\begin{array}{c c c c c c c c}
f(1) & f(2) & f(3) & f(4) & f(5) & f(6) & f(7) & \dots \\ \\
- \lceil \tfrac{0}{2} \rceil & \lceil \tfrac{1}{2}\rceil & - \lceil \tfrac{2}{2} \rceil & \lceil \tfrac{3}{2} \rceil & - \lceil \tfrac{4}{2} \rceil  & \lceil \tfrac{5}{2}
\rceil & - \lceil \tfrac{6}{2} \rceil & \dots \\ \\
0 & 1 & -1 & 2 & -2 & 3 & \dots
\end{array}
\]
Kita uga bisa nganggep $f$ ditetepake miturut kasus kaya mangkene:
\[
f(n) = \begin{cases}
  0 & \text{yen $n = 1$}\\
  n/2 & \text{yen $n$ genep}\\
  -(n-1)/2 & \text{yen $n$ ganjil lan $>1$}
  \end{cases}
\]
\end{ex}

\begin{prob}
  Tuduhna yen $A$ lan $B$ iku !!{enumerable}, mula $A \cup B$ uga mangkono.
  Kanggo iku, upamakna ana fungsi !!{surjective} $f\colon \PosInt \to
  A$ lan $g\colon \PosInt \to B$, banjur tetepna fungsi !!a{surjective}
  $h\colon \PosInt \to A \cup B$ lan buktekna yen fungsi iku
  !!{surjective}. Uga gatekna kasus nalika $A$ utawa~$B = \emptyset$.
\end{prob}

\begin{prob}
  Tuduhna yen $B \subseteq A$ lan $A$ iku !!{enumerable}, mula~$B$ uga mangkono.
  Kanggo iku, upamakna ana fungsi !!a{surjective} $f\colon \PosInt \to
  A$. Tetepna fungsi !!a{surjective} $g\colon \PosInt \to B$ lan buktekna
  yen fungsi iku !!{surjective}. Apa kang kedadeyan yen $B = \emptyset$?
\end{prob}

\begin{prob}
  Kanthi induksi tumrap $n$, tuduhna yen $A_1$, $A_2$, \dots, $A_n$ kabeh
  !!{enumerable}, mula $A_1 \cup \dots \cup A_n$ uga mangkono. Kita oleh nggunakake
  asil yen rong himpunan $A$ lan~$B$ iku !!{enumerable}, mula~$A \cup
  B$ uga mangkono.
\end{prob}

Sanajan nalika ndhaptar !!{element}s himpunan katon luwih lumrah
miwiti saka !!{element} kaping $1$, ahli matematika seneng
nggunakake wilangan natural~$\Nat$ kanggo ngetung.
Dheweke nyebut !!{element}s dhaptar kaping $0$, kaping $1$, kaping $2$, lan sateruse.
Kanthi mangkono, kita bisa netepake enumerasi minangka
fungsi !!a{surjective} saka $\Nat$ menyang~$A$.
Mesthi wae, definisi loro iku ekuivalen.

\begin{prop}\ollabel{prop:enum-shift}
  Ana !!a{surjection} $f\colon \PosInt \to A$ yen lan mung yen ana
  !!a{surjection} $g\colon \Nat \to A$.
\end{prop}

\begin{proof}
  Yen diwenehi !!a{surjection} $f\colon \PosInt \to A$, kita bisa netepake $g(n) =
  f(n+1)$ kanggo kabeh $n \in \Nat$. Gampang dideleng yen $g\colon \Nat
  \to A$ iku !!{surjective}. Kosok baline, yen diwenehi !!a{surjection} $g\colon
  \Nat \to A$, tetepna $f(n) = g(n-1)$.
\end{proof}

Saka iki kita oleh asil:

\begin{cor}\ollabel{cor:enum-nat}
Himpunan $A$ iku !!{enumerable} yen lan mung yen kosong utawa ana
fungsi !!a{surjective} $f\colon \Nat \to A$.
\end{cor}

Ing ndhuwur kita wis ngrembug yen dhaptar !!{element}s himpunan~$A$ bisa
diowahi dadi dhaptar tanpa pangulangan. Iki uga lumaku tumrap
enumerasi, nanging rada luwih angel dirumusake lan dibuktekake kanthi tliti.
Saben fungsi $f\colon \PosInt \to A$ kudu ditetepake kanggo kabeh $n \in
\PosInt$. Yen cacah !!{element}s ing~$A$ mung winates, cetha yen ora bisa
ana fungsi kang ditetepake ing !!{element}s saka~$\PosInt$ kang cacahe tanpa wates,
nggayuh kabeh !!{element}s saka~$A$ minangka nilai, nanging ora nate
nggayuh nilai kang padha kaping pindho. Ing kasus iku, yaiku
nalika dhaptar tanpa pangulangan winates, kita kudu milih domain
liyane kanggo~$f$, kang cacah !!{element}s mung winates.
Tanpa pangulangan ateges $f$ kudu !!{injective}. Amarga uga
!!{surjective}, kita nggoleki !!a{bijection} antarane sawijining
himpunan winates $\{1, \dots, n\}$ utawa $\PosInt$ lan~$A$.

\begin{prop}\ollabel{prop:enum-bij}
Yen $f\colon \PosInt \to A$ iku !!{surjective}, yaiku enumerasi
saka~$A$, ana !!a{bijection} $g\colon Z \to A$ kanthi $Z$
yaiku~$\PosInt$ utawa $\{1, \dots, n\}$ kanggo sawijining~$n \in \PosInt$.
\end{prop}

\begin{proof}
  Kita netepake fungsi $g$ kanthi rekursif: tetepna $g(1) = f(1)$. Yen $g(i)$
  wis ditetepake, tetepna $g(i+1)$ minangka nilai kapisan saka $f(1)$,
  $f(2)$, \dots{} kang durung ana ing antarane $g(1)$, \dots, $g(i)$, yen ana.
  Yen $A$ mung duwe $n$ !!{element}s, mula $g(1)$, \dots, $g(n)$ kabeh
  ditetepake, saengga kita wis netepake fungsi $g\colon \{1, \dots, n\}
  \to A$. Yen cacah !!{element}s ing $A$ tanpa wates, mula kanggo saben $i$
  mesthi ana !!a{element} saka~$A$ ing enumerasi $f(1)$, $f(2)$,
  \dots, kang durung ana ing antarane $g(1)$, \dots, $g(i)$. Ing kasus
  iki kita wis netepake fungsi $g\colon \PosInt \to A$.

  Fungsi $g$ iku !!{surjective}, amarga saben anggota saka~$A$ ana
  ing antarane $f(1)$, $f(2)$, \dots{} (amarga $f$ iku !!{surjective}) lan
  pungkasane bakal dadi nilai~$g(i)$ kanggo sawijining~$i$. Fungsi iki uga
  !!{injective}, amarga yen ana $j < i$ saengga $g(j) = g(i)$,
  mula $g(i)$ wis ana ing antarane $g(1)$, \dots, $g(i-1)$, kang
  nalisir cara kita netepake~$g$.
\end{proof}

\begin{cor}\ollabel{cor:enum-nat-bij}
Himpunan $A$ iku !!{enumerable} yen lan mung yen kosong utawa ana !!a{bijection}
$f\colon N \to A$ kanthi $N = \Nat$ utawa $N = \{0, \dots, n\}$ kanggo
sawijining $n \in \Nat$.
\end{cor}

\begin{proof}
$A$ iku !!{enumerable} yen lan mung yen $A$ kosong utawa ana fungsi !!a{surjective}
$f\colon \PosInt \to A$. Miturut \olref{prop:enum-bij}, sarat pungkasan iku
lumaku yen lan mung yen ana fungsi !!a{bijective}~$f\colon Z \to A$ kanthi $Z =
\PosInt$ utawa $Z = \{1, \dots, n\}$ kanggo sawijining $n \in \PosInt$. Kanthi
argumen kang padha karo bukti \olref{prop:enum-shift}, sarat iku sabanjure
lumaku yen lan mung yen ana !!a{bijection} $g\colon N \to A$ kanthi
$N = \Nat$ utawa $N = \{0, \dots, n-1\}$.
\end{proof}

\begin{prob}
  Miturut \olref[sfr][siz][enm]{defn:enumerable}, himpunan $A$
  bisa didhaptar yen lan mung yen $A = \emptyset$ utawa ana fungsi !!a{surjective} $f\colon
  \PosInt \to A$. Uga bisa netepake "himpunan kang !!{enumerable}"
  kanthi tliti: himpunan bisa didhaptar yen lan mung yen ana fungsi !!a{injective}
  $g\colon A \to \PosInt$. Tuduhna yen definisi loro iku
  ekuivalen, yaiku ana fungsi !!a{injective}
  $g\colon A \to \PosInt$ yen lan mung yen $A = \emptyset$ utawa ana
  fungsi !!a{surjective} $f\colon \PosInt \to A$.
\end{prob}

\end{document}
